\documentclass[11pt]{article}
\usepackage[margin=1in]{geometry}
\usepackage{booktabs}
\usepackage{hyperref}
\usepackage{microtype}
\usepackage{amsmath}
\hypersetup{colorlinks=true, linkcolor=blue, urlcolor=blue, citecolor=blue}

\title{Splendor as a Fast Reinforcement Learning Environment:\\
League Self-Play, Elo Evaluation and Search}
\author{Aditya Kumar\\ \small Carnegie Mellon University \quad \texttt{adityaku@andrew.cmu.edu}}
\date{September 2026}

\begin{document}
\maketitle

\begin{abstract}
We present an exact-rules implementation of the board game \emph{Splendor} (2--4 players) as a
reinforcement learning environment for PufferLib. The game logic is a single C header that
advances roughly four million turns per second on one CPU core, exposes a legal-action mask
inside the observation, and supports state snapshots for tree search. On top of it we build a
prioritized fictitious self-play league for PPO, a Bradley--Terry Elo ladder with fixed
heuristic yardsticks, a batched PUCT Monte Carlo tree search with an AlphaZero-style training
loop, and a browser interface for watching or playing against agents. A 300M-step league run
on a single workstation produces a 2-player policy rated about 230 Elo above a strong greedy
heuristic and 215 Elo above plain self-play with the same network and budget; 64 simulations of
search add a further 100 Elo. Code, trained models and the interface are released under the MIT
license at \url{https://github.com/Sylindril/splendor}.
\end{abstract}

\section{Introduction}
Splendor is a resource-management card game with hidden information (deck order and
face-down reserved cards), chance (card draws), sub-turn decisions (returning tokens above the
ten-token limit, choosing among nobles) and two to four players. These properties make it a
useful mid-sized testbed between fully observable two-player board games and large
multi-agent simulations. Existing implementations are written for human play or for slow
Python search and are not suitable for training policies with hundreds of millions of samples.

This note describes a from-scratch environment designed for throughput and simplicity, and the
training and evaluation tooling around it. The contributions are practical: (i) an exact-rules
engine with a compact fixed-size observation and a 72-way discrete action space that is
verified by sanitizer runs and invariant tests; (ii) a league training wrapper that removes
idle transitions and mixes the current policy with a pool of past checkpoints sampled by
prioritized fictitious self-play (PFSP)~\cite{vinyals2019}; (iii) an Elo ladder built on
Bradley--Terry ratings~\cite{hunter2004} with fixed bots as anchors; (iv) a batched PUCT search
that uses engine snapshots and determinization, and (v) shipped models with an interface.

\section{Environment}
\paragraph{Rules.} The engine implements the published rules without simplification: 90
development cards in three tiers with four face-up per tier, ten nobles of which $P+1$ are
dealt, a bank of 4/5/7 gems per color for 2/3/4 players plus five gold jokers, the four turn
actions (take up to three different gems, take two of one color when at least four remain,
reserve a face-up or deck-top card gaining a gold, buy a card paying with bonuses then gems
then gold), the ten-token limit enforced by a discard phase, noble visits after a purchase
with a choice when several qualify, and game end at the end of the round in which a player
reaches 15 points, with the fewest-cards tie-break. The start seat is randomized per game so
that seat index carries no advantage. The card table was cross-checked against four
independent open-source implementations.

\paragraph{Interface.} One C struct holds one game; each seat is a PufferLib agent, so a
vectorized environment of $N$ games exposes $NP$ agents. One environment step is one turn of
the seat to move; other seats' actions are ignored and their observations carry a pass-only
mask. Sub-phases (discard, noble choice) are additional steps by the same seat. The
observation is a fixed vector of $240 + 48P$ bytes: bank, deck sizes, twelve face-up cards
(cost, bonus, points), five noble requirement vectors, per-player blocks (tokens, bonuses,
points, three reserved cards with a presence flag; opponents' deck-top reserves are hidden),
the 72-way legal-action mask, the turn number and a to-move flag. Policies read the mask from
the observation and mask their logits, so no invalid action is ever sampled and the trainer
needs no masking support. The engine additionally provides opaque state snapshots, restore,
a no-op action that freezes a game, and a \emph{determinize} operation that re-deals
everything a given seat cannot see.

\paragraph{Throughput.} On an Apple M4 Max, a single core advances 4.1M turns/s with random
actions and 1.9M turns/s when actions are sampled uniformly from the mask in NumPy. Address
and undefined-behavior sanitizers over one million random steps at each player count report
no findings; invariant tests check token, card and noble conservation, the ten-token limit,
and that illegal actions are no-ops.

\section{Training}
\paragraph{Policy.} A three-layer, 512-unit LayerNorm MLP over the scaled observation with
masked actor and value heads (0.74M parameters). Plain self-play, in which every seat is the
current policy, works but wastes the idle seats' transitions and can cycle.

\paragraph{League.} A wrapper exposes only the learner's seat to the trainer. Inside one
wrapper step it applies the learner's action, then plays every opponent seat (including their
sub-phases) with frozen policies evaluated on CPU inside the worker process, batched per
model, until the learner is to move again. Every trainer sample is therefore a real decision.
Opponents are assigned per game: with probability $\tfrac12$ the latest weights (reloaded
from disk every few seconds), otherwise a member of a pool of past checkpoints sampled with
PFSP weights $(1-w_i)^2 + 0.05$, where $w_i$ is an exponential moving average of the learner's
win rate against member $i$. A checkpoint joins the pool every 20M steps. The PPO
implementation is PufferLib's~\cite{suarez2024}, with V-trace corrections and prioritized
minibatches, reward $\pm1$ for win/loss plus $0.02$ per prestige point. The wrapper sustains
about 180K learner steps per second per process on the reference machine; end-to-end training
runs at 150--180K samples per second with eight workers and an MPS device.

\section{Evaluation}
\paragraph{Ladder.} Ratings are fitted with the Bradley--Terry model by
minorization--maximization on pairwise results (draws count half), with a one-game prior so
that undefeated participants remain finite, reported on the Elo scale ($400\log_{10}$) and
anchored at a uniform-random-legal agent. Two-player ladders are round robins; three- and
four-player ladders sample random tables and extract pairwise results from final ranks. A
fixed greedy heuristic (buy the most valuable affordable card, else take gems toward the
closest card) serves as a stable yardstick: it beats random in 100\% of games.

\begin{table}[h]
\centering
\begin{tabular}{lrr}
\toprule
Participant (2 players, 100 games per pair) & Elo & Win \% \\
\midrule
League, 300M steps, final & 1210 & 75.7 \\
League pool, 240M & 1197 & 74.5 \\
League pool, 160M & 1049 & 58.0 \\
Plain self-play, same network, 60M & 994 & 52.4 \\
Greedy heuristic & 982 & 51.0 \\
League pool, 80M & 945 & 46.8 \\
League pool, 40M & 740 & 28.9 \\
Random & 0 & 6.7 \\
\bottomrule
\end{tabular}
\caption{Two-player ladder. The league policy gains about 230 Elo over the greedy bot and 215
over plain self-play; the curve is flattening but still rising at 300M steps.}
\end{table}

\begin{table}[h]
\centering
\begin{tabular}{lr}
\toprule
Participant (4 players, 150 random tables) & Elo \\
\midrule
League, 200M steps, final & 1017 \\
League pool, 175M & 1009 \\
League pool, 125M & 971 \\
Greedy heuristic & 894 \\
League pool, 75M & 886 \\
League pool, 25M & 699 \\
Random & 0 \\
\bottomrule
\end{tabular}
\caption{Four-player ladder after a shorter run; the margin over the heuristic is smaller
(120 Elo) with a third of the per-seat experience and higher outcome variance.}
\end{table}

\section{Search}
The search is a batched PUCT over engine snapshots: per simulation round every root descends
to a leaf in lock-step, one vectorized environment step expands all leaves, and one network
forward pass evaluates them. Values are per-seat vectors (the network is evaluated on every
seat's observation of the leaf), so multiplayer backups are exact; terminal leaves use the
engine's rewards. Hidden information is handled by determinizing the root for the seat to
move, and the engine's random state is part of the snapshot, so a search explores one
determinization. With the 512-unit network the search runs about 26K simulations per second
on CPU at batch 64. Sixty-four simulations per move beat the raw policy they search with in
63\% of games, which the ladder reports as roughly 100 Elo. An AlphaZero-style loop (self-play
with search, then cross-entropy on visit distributions plus value regression) is included and
warm-starts from a PPO checkpoint; scaling it is left to future work.

\section{Interface}
A dependency-free browser interface renders a square table with a mat per player, animates
every move from a structured event emitted by the server (cards flying to mats, replacement
cards flipping in from the deck, gems moving to and from the bank), holds the finished board
for inspection with standings, supports keyboard control, and lets the user seat any mix of
humans, bots, checkpoints and search agents.

\section{Limitations}
No human rating is available, so ``superhuman'' claims cannot be made; the ladder measures
strength relative to bots and past checkpoints only. Rewards include small prestige-point
shaping. The search uses a single determinization per move. Runs used one workstation; the
league is designed to scale to many workers on a GPU host but this was not benchmarked here.

\section{Conclusion}
A small, exact and fast environment plus a league wrapper gives a clear and cheap path from
random play to a policy well beyond a strong heuristic in under an hour on a desktop, with
search as a further multiplier. We hope the code and models are useful as a reproducible
mid-sized benchmark for multi-agent reinforcement learning.

\begin{thebibliography}{9}
\bibitem{suarez2024} J.~Suarez. PufferLib: Making reinforcement learning libraries and
environments play nice. \emph{arXiv:2406.09411}, 2024.
\bibitem{vinyals2019} O.~Vinyals et al. Grandmaster level in StarCraft II using multi-agent
reinforcement learning. \emph{Nature}, 575:350--354, 2019.
\bibitem{hunter2004} D.~R.~Hunter. MM algorithms for generalized Bradley--Terry models.
\emph{Annals of Statistics}, 32(1):384--406, 2004.
\bibitem{silver2018} D.~Silver et al. A general reinforcement learning algorithm that masters
chess, shogi, and Go through self-play. \emph{Science}, 362(6419):1140--1144, 2018.
\end{thebibliography}
\end{document}
